% !PRINTED 563
\weberpage{563}
%<<CONT
\[
\xi = r + r^{-1} = 2\cos\frac{2\pi}{13}, \qquad \xi' = r^5 + r^{-5} = 2\cos\frac{10\pi}{13}
\]

sets. For we have

\[
\xi + \xi' = \eta, \qquad \xi\xi' = \eta_2,
\]

hence

\weq{11}{\xi^2 - \eta\xi + \eta_2 = 0}

is the quadratic equation for $\xi$, from which one obtains the biquadratic
equation for $r$,

\weq{12}{r^4 - \eta r^3 + (\eta_2 + 2)r^2 - \eta r + 1 = 0.}

The quadratic equation (11) has one positive and one negative
root. The positive root is $2\cos\frac{2\pi}{13}$. Instead of computing $r$
itself, one computes further $2\sin\frac{2\pi}{13}$, which one obtains as
the positive square root $\sqrt{4 - \xi^2}$.

In place of the three four-term periods $\eta$ one could also
have first considered two periods of six terms.

\weq{13}{\begin{aligned}
\xi &= r + r^3 + r^4 + r^{-1} + r^{-3} + r^{-4}\\
\xi_1 &= r^2 + r^5 + r^6 + r^{-2} + r^{-5} + r^{-6}.
\end{aligned}}

For the product $\xi\xi_1$ one finds by formula (4) the
value $3(\xi + \xi_1)$ or $-3$, so that $\xi$, $\xi_1$ are the roots of the quadratic equation

\weq{14}{\xi^2 + \xi - 3 = 0}

from which, since it is easy to see from the expression by the
cosines that $\xi$ is positive:

\weq{15}{\xi = \frac{-1 + \sqrt{13}}{2}, \qquad \xi_1 = \frac{-1 - \sqrt{13}}{2}.}

After adjunction of the values $\xi$, $\xi_1$ one can form a cubic
equation for the two-term periods. For if one
sets

\[
\begin{aligned}
\xi &= r + r^{-1}, & \xi_1 &= r^2 + r^{-2},\\
\xi_2 &= r^4 + r^{-4}, & \xi_3 &= r^5 + r^{-5},\\
\xi_4 &= r^3 + r^{-3}, & \xi_5 &= r^6 + r^{-6},
\end{aligned}
\]

then it follows that

\[
\zeta = \xi + \xi_2 + \xi_4, \qquad \zeta_1 = \xi_1 + \xi_3 + \xi_5, \qquad \zeta + \zeta_1 = -1,
\]

further

\[
\xi\xi_2 = \xi_3 + \xi_4, \qquad \xi\xi_4 = \xi_1 + \xi_2, \qquad \xi_2\xi_4 = \xi + \xi_5,
\]

\[
\xi\xi_2\xi_4 = 2 + \zeta_1 = 1 - \zeta,
\]
%CONT>>
